\section{Discussion}\label{sec:discussion}
\subsection{How the approach answers the open problems}
\runin{G1, evaluation.} Every reported number comes from held-out photos or videos, with a false-alarm metric, a second wild photo set and five standard baselines on the same rows. The honest outcome is mixed: the cascade removes most false alarms of the rule engine (78.3\% to 20.8\%) and holds up on wild photos, but a random forest is stronger on photos from the training sources, which a random-split study would not have revealed.
\runin{G2, temporal behaviour.} The flow check recognises held poses and separates them from transitions on videos of unseen people (78.0\% macro recall, three poses above the bar) and its input-rate requirement is measured. A controlled comparison shows that simpler sequence models are as good for held-pose recognition and at least as good on a rule-defined movement task, and a time-reversal test is at chance for every model, so neither the graph nor the frame order is shown to be necessary, whereas the depth coordinate that only the sequence input carries is shown to matter (removing it costs every sequence model 4 to 60 points); our labels support only one named transition, so movement modelling is untested, and the flow check does not yet score smoothness, balance or momentum.
\runin{G3, personalisation.} A 15-second, vision-only range profile adapts the score without extra sensors; the mechanism behaves as designed but has not been validated with users.
\runin{G4, hidden joints.} The system declines to score what it cannot see, and the measured alternatives show that mirroring is better for symmetric stances and mesh recovery for legs and asymmetric poses, but the latter needs the image.
\runin{G5, generated advice.} Cues come from reviewed templates, an optional paraphrase may only rewrite them, a term screen can discard the rewrite, and repeated failures back off instead of pushing harder.
\runin{G6, deployment.} Only landmark numbers leave the device, no extra hardware is needed, guidance is available in three languages, and a rule-engine port keeps a basic mode working offline.

\subsection{Cost-efficiency}
Recent work often raises model complexity for marginal accuracy gains. Our design spends parameters sparingly: two 0.36\,M-parameter MLPs and one 0.89\,M-parameter ST-GCN run in 324\,MB on a 2-vCPU host at 7--9\,ms per frame. The measurements show where complexity pays and where it does not. The second, small MLP (the gate) costs about one more forward pass and cuts false alarms from 81.5\% (namer alone) to 20.8\%; residual blocks add nothing over a plain MLP on public photos, and an 86\,k-parameter MLP on window statistics matches or exceeds the 0.89\,M-parameter ST-GCN at held-pose recognition while training in 6\,s instead of 290\,s per fold; and mesh recovery improves the median angle error by $2.8^\circ$ at 164\,ms per image on a CPU and needs the camera image, so we do not use it.

\subsection{Limitations}
\begin{itemize}
\item \emph{No live-person accuracy study.} All end-to-end tests use photographs or recorded video; a study with practitioners and instructors is needed.
\item \emph{The form score is not validated against human judgement.} It is checked only by a perturbation test, it is nearly blind to arm errors (0.94 to 0.92 when an elbow is broken), and the learned deviation head is close to chance (12.4\% against 6.7\%); joint-level cues rest on hand-set angle bands, and Tree and Lunge have none.
\item \emph{Weak poses.} Mountain, cobra and child's pose fail the phone photo test; the flow check reliably handles three poses and is matched by far simpler sequence models; our labels name only one transition (cobra pose to downward dog) with support in three or more videos, so movement modelling is untested; named transitions and momentum scores are not implemented.
\item \emph{Small data.} 24 videos and public photo sets; the held-out public photos share their sources with the training photos; the wild set has 103 photos; baselines use one seed.
\item \emph{Occlusion.} Pose naming can still be influenced by a hallucinated hidden joint; the mesh-recovery study is small (141 cases, one asymmetric pose, a synthetic occluder).
\item \emph{Personal profile and safety.} The profile is not user-tested and a poor stance could forgive real errors; a term screen cannot certify open-ended text, the Hindi and Bengali templates and term lists were checked by one native speaker of both languages on the team and have not had an independent review by instructors.
\item \emph{Latency and hardware.} The round trip is about 1\,s from India because it is network-dominated; one phone model was tested. The system is not a medical device and no injury outcome was studied.
\end{itemize}

\subsection{Future work}
(i)~A live user study in which instructors rate form, producing the labelled good/bad-form data the form head lacks, in the manner of recent isometric-pose benchmarks \cite{jaiswal2025isometric}; (ii)~a tree--network ensemble for naming, since the two families win on different sources; (iii)~on-device inference of the 0.36\,M-parameter MLPs, which would remove the round trip; (iv)~hand-set bands for Tree and Lunge; (v)~flow scores (smoothness, balance, momentum) and named transitions, the tasks on which a temporal model must prove its value over static statistics; (vi)~retraining the form head with arm perturbations; (vii)~viewpoint-robust features and a larger pose vocabulary; (viii)~opt-in server-side mesh recovery with explicit image consent.

\section{Conclusion}
AsanaAI shows that a yoga coach can keep the camera image on the device, abstain when evidence is missing and still run on commodity hardware, provided it is evaluated on data it never trained on. A name-and-veto cascade of small residual MLPs names poses on 78.9\% of held-out photos with 20.8\% false alarms, a skeleton-graph flow check reaches 78.0\% macro recall on held-out videos although simpler sequence models are as good on held poses and on a rule-defined movement task, and the system degrades gracefully through three pose engines and an offline mode. The comparison with standard methods shows that no single classifier wins on every photo source, and we report what the system cannot yet do.

\runin{Availability.} App: \url{https://yoga-posture-correction-system.vercel.app}; code: \url{https://github.com/Anamitra-Sarkar/yoga-posture-correction-system}; training data and code: Hugging Face dataset Arko007/Yoga-1M (transcripts withheld because they are copyrighted speech).
